\ifx\allfiles\undefined
\documentclass[12pt, a4paper, oneside, UTF8]{ctexbook}
\def\path{../config}
\input{\path/package.tex}

% 定理环境设置
\input{\path/theorem1_zh.tex}

\input{\path/custom.tex}

% 修改参数改变封面样式，0 默认原始封面、内置其他1、2、3种封面样式
\def\myIndex{3}


\ifnum\myIndex>0
    \input{\path/cover_package_\myIndex} 
\fi

\def\myTitle{线性代数笔记}
\def\myAuthor{M.K.}
\def\myDateCover{\today}
\def\myDateForeword{\today}
\def\myForeword{前言}
\def\myForewordText{
    
    这是一个基于\LaTeX{}模板的线代笔记。

}
\def\mySubheading{副标题}


\begin{document}
% \input{\path/cover_text_\myIndex.tex}

\newpage
\thispagestyle{empty}
\begin{center}
    \Huge\textbf{\myForeword}
\end{center}
\myForewordText
\begin{flushright}
    \begin{tabular}{c}
        \myDateForeword
    \end{tabular}
\end{flushright}

\newpage
\pagestyle{plain}
\setcounter{page}{1}
\pagenumbering{Roman}
\pagestyle{fancy}
\fancyfoot[C]{\thepage}
\renewcommand{\headrulewidth}{0.4pt}
\renewcommand{\footrulewidth}{0pt}
\tableofcontents

\newpage
\pagenumbering{arabic}
\setcounter{page}{1}








\else
\fi
\chapter{二次型}

\section{二次型矩阵}

\begin{definition}[二次型]
    对于含有$n$个变量$x_1, x_2, \cdots, x_n$的二次齐次函数
    \begin{align*}
        f(x_1, x_2, \cdots, x_n) =& \quad\,a_{11} x_1^2 + a_{22} x_2^2 + \cdots + a_{nn} x_n^2 \\
        &+ 2 a_{12} x_1 x_2 + 2 a_{13} x_1 x_3 + \cdots + 2 a_{1n} x_1 x_n\\
        &+ 2 a_{23} x_2 x_3 + \cdots + 2 a_{2n} x_2 x_n \\
        &+ \cdots \\
        &+ 2 a_{n-1,n} x_{n-1} x_n
    \end{align*}
    存在一个唯一的对称矩阵$A = (a_{ij})_{n \times n}$，使得
    \[ f(x_1, x_2, \cdots, x_n) = \bm{x}^T A \bm{x} \]
    其中$\bm{x} = (x_1, x_2, \cdots, x_n)^T$．称$A$为二次型$f$的对应矩阵．
    \begin{myRemark}
        若题目给出的$\bm{A}$不是对称矩阵，则应先将其化为对称矩阵$\frac{\bm{A} + \bm{A}^T}{2}$，再作为对应矩阵．
    \end{myRemark}
\end{definition}

\begin{theorem}
    对于任意给定的二次型，其对应的对称矩阵是唯一的．
\end{theorem}

\section{二次型的变换}

\subsection{二次型与线性变换}

\begin{definition}[非退化线性变换]
    若线性变换 $\bm{x} = \bm{C}\bm{y}$ 的系数矩阵 $\bm{C}$ 可逆，则称为非退化线性变换．\\
    非退化线性变换把二次型 $f = \bm{x}^T \bm{A} \bm{x}$ 变为 $g = \bm{y}^T (\bm{C}^T \bm{A} \bm{C}) \bm{y}$，矩阵 $\bm{C}^T \bm{A} \bm{C}$ 与 $\bm{A}$ 合同．
\end{definition}

\begin{definition}[标准型]
    二次型经非退化线性变换可化为如下形式，称为标准型：
    \[ f = k_1 y_1^2 + k_2 y_2^2 + \cdots + k_n y_n^2 \]
\end{definition}

\begin{definition}[惯性指数]
    在二次型的标准型中，正系数的个数 $p$ 称为正惯性指数，负系数的个数 $q$ 称为负惯性指数．
\end{definition}

\begin{definition}[规范型]
    二次型经非退化线性变换可化为如下形式，称为规范型：
    \[ f = y_1^2 + \cdots + y_p^2 - y_{p+1}^2 - \cdots - y_{p+q}^2 \]
    其中 $p$ 为正惯性指数，$q$ 为负惯性指数．
\end{definition}

\begin{theorem}[惯性定理]
    实二次型的标准型虽不唯一，但其规范型是唯一的．即正惯性指数 $p$ 和负惯性指数 $q$ 由二次型本身唯一确定．二次型对应的矩阵为$\bm{A}$，则有
    \[p + q = \rank(\bm{A})\]
\end{theorem}

\subsection{矩阵合同}

\begin{definition}[合同]
    设 $\bm{A}, \bm{B}$ 为 $n$ 阶矩阵，若存在可逆矩阵 $\bm{C}$，使得 $\bm{C}^T \bm{A} \bm{C} = \bm{B}$，则称矩阵 $\bm{A}$ 与 $\bm{B}$ 合同，记作 $\bm{A} \simeq \bm{B}$．
\end{definition}

\begin{theorem}
    实对称矩阵 $\bm{A}$ 与 $\bm{B}$ 合同的充要条件是它们有相同的正惯性指数 $p$ 和负惯性指数 $q$（即具有相同的惯性指数）．
\end{theorem}

\begin{theorem}
    实对称矩阵只能和实对称矩阵合同．\\
    $\text{\qquad\qquad\qquad\quad}$非实对称只能和非实对称矩阵合同．
\end{theorem}
\subsection{配方法}

\begin{enumerate}
    \item 若二次型含平方项 $x_i^2$（不妨设 $a_{ii} \neq 0$），则选取含 $x_i$ 的所有项进行配方．
    \item 若二次型不含平方项，但含交叉项 $a_{ij} x_i x_j$，则通过变换
    \[ \begin{cases} x_i = y_i - y_j \\ x_j = y_i + y_j \end{cases} \]
    引入平方项，再按第一种情况处理．
\end{enumerate}

\subsection{初等变换法}

对对称矩阵 $\bm{A}$ 作初等行变换，同时作对应的初等列变换，将其化为对角矩阵．
\[ \bm{A} \xrightarrow{r_i + k r_j} \bm{A}_1 \xrightarrow{c_i + k c_j} \bm{A}_2 \xrightarrow{\cdots} \bm{\Lambda} \]
若需求变换矩阵 $\bm{C}$，可构造 $(\bm{A} \mid \bm{E})$，对列作变换的同时作用于 $\bm{E}$，对行作变换时不作用于 $\bm{E}$．最后右侧的 $\bm{E}$ 变为 $\bm{C}$．

\subsection{正交变换法}

利用正交矩阵 $\bm{Q}$ 将 $\bm{A}$ 对角化（$\bm{Q}^T \bm{A} \bm{Q} = \bm{\Lambda}$）：
\begin{enumerate}
    \item 求出 $\bm{A}$ 的所有特征值 $\lambda_1, \cdots, \lambda_n$．
    \item 求出对应特征向量，并进行正交化、单位化．
    \item 构造正交矩阵 $\bm{Q}$，作变换 $\bm{x} = \bm{Q}\bm{y}$，得到标准型 $f = \lambda_1 y_1^2 + \cdots + \lambda_n y_n^2$．
\end{enumerate}

\begin{example}
    设二次型$f(x_1, x_2, x_3) = \displaystyle\sum_{i=1}^{3}\sum_{j=1}^{3}ijx_ix_j$
    \begin{tasks}[label=(\arabic*)]
        \task 利用正交变换$\bm{x}=\bm{Qy}$将$f(x_1, x_2, x_3)$化为标准型；
        \task 求$f(x_1, x_2, x_3)=0$的解.
    \end{tasks}
\end{example}
\begin{solution}
    二次型对应的矩阵为
    \[ \bm{A} = \begin{bNiceMatrix}[columns-width = 1em] 1 & 2 & 3 \\ 2 & 4 & 6 \\ 3 & 6 & 9 \end{bNiceMatrix} = \begin{bNiceMatrix} 1 \\ 2 \\ 3 \end{bNiceMatrix} \begin{bNiceMatrix} 1 & 2 & 3 \end{bNiceMatrix} \]
    显然$\rank(\bm{A})=1$，特征值为$\lambda_1=14, \lambda_2=\lambda_3=0$.
    \begin{enumerate}
        \item 对应$\lambda_1=14$的特征向量$\bm{\xi}_1 = (1, 2, 3)^T$，单位化得$\bm{\eta}_1 = \frac{1}{\sqrt{14}}(1, 2, 3)^T$.

        对应$\lambda_2=\lambda_3=0$的特征向量需满足$x_1+2x_2+3x_3=0$，取$\bm{\xi}_2 = (2, -1, 0)^T$，$\bm{\xi}_3$需与$\bm{\xi}_1, \bm{\xi}_2$正交，可取$\bm{\xi}_1 \times \bm{\xi}_2 = (3, 6, -5)^T$.

        单位化得$\bm{\eta}_2 = \frac{1}{\sqrt{5}}(2, -1, 0)^T, \bm{\eta}_3 = \frac{1}{\sqrt{70}}(3, 6, -5)^T$.

        令$\bm{Q} = (\bm{\eta}_1, \bm{\eta}_2, \bm{\eta}_3)$，则经$\bm{x}=\bm{Qy}$变换得标准型
        \[ f = 14y_1^2 \]
        \item $f=0 \implies 14y_1^2=0 \implies y_1=0$.

        由$\bm{y}=\bm{Q}^T\bm{x}$知$\bm{x}=\bm{Qy}=\bm{Q}\begin{bNiceMatrix} 0 \\ y_2 \\ y_3 \end{bNiceMatrix} = y_2 \bm{\eta}_2 + y_3 \bm{\eta}_3$，即
        \[ \bm{x} = k_1 \begin{bNiceMatrix} 2 \\ -1 \\ 0 \end{bNiceMatrix} + k_2  \begin{bNiceMatrix} 3 \\ 6 \\ -5 \end{bNiceMatrix} \]
        其中$k_1, k_2$为任意实数.
    \end{enumerate}
\end{solution}

\begin{example}
    设二次型$f(x_1, x_2) = x_1^2 - 4x_1 x_2 + 4x_2^2$经正交变换$\begin{bNiceMatrix} x_1 \\ x_2 \end{bNiceMatrix} = \bm{Q} \begin{bNiceMatrix} y_1 \\ y_2 \end{bNiceMatrix}$化为一个新的二次型$g(y_1, y_2) = ay_1^2+4y_1y_2+by_2^2$，其中$a \ge b$，:\begin{tasks}[label=(\arabic*)](2)
        \task 求常数$a, b$;
        \task 给出一个正交矩阵$\bm{Q}$.
    \end{tasks}
\end{example}
\begin{solution}
    \begin{enumerate}
        \item 二次型$f, g$的矩阵分别为$\bm{A} = \begin{bNiceMatrix} 1 & -2 \\ -2 & 4 \end{bNiceMatrix}$, $\bm{B} = \begin{bNiceMatrix}[columns-width = 1.2em] a & 2 \\ 2 & b \end{bNiceMatrix}$.

        因$f, g$经正交变换相互转化，故$\bm{A}, \bm{B}$相似，有相同的特征值，从而迹与行列式相等：
        \[ \begin{cases} \operatorname{tr}(\bm{A}) = 1+4 = a+b = \operatorname{tr}(\bm{B}) \\ \det(\bm{A}) = 4-4 = ab-4 = \det(\bm{B}) \end{cases} \implies \begin{cases} a+b=5 \\ ab=4 \end{cases} \]
        解得$\begin{cases} a=1 \\ b=4 \end{cases}$或$\begin{cases} a=4 \\ b=1 \end{cases}$，因$a \ge b$，故$a=4, b=1$.
        \item  $\bm{A} = \begin{bNiceMatrix} 1 & -2 \\ -2 & 4 \end{bNiceMatrix}$的特征值为$\lambda_1=5, \lambda_2=0$.

        对应$\lambda_1=5$，$\bm{\xi}_1 = \frac{1}{\sqrt{5}}(1, -2)^T$；对应$\lambda_2=0$，$\bm{\xi}_2 = \frac{1}{\sqrt{5}}(2, 1)^T$.

        令$\bm{Q}_1 = (\bm{\xi}_1, \bm{\xi}_2) = \frac{1}{\sqrt{5}}\begin{bNiceMatrix} 1 & 2 \\ -2 & 1 \end{bNiceMatrix}$.

        $\bm{B} = \begin{bNiceMatrix} 4 & 2 \\ 2 & 1 \end{bNiceMatrix}$同理可得特征值$\lambda_1=5, \lambda_2=0$.

        对应$\lambda_1=5$，$\bm{\eta}_1 = \frac{1}{\sqrt{5}}(2, 1)^T$；对应$\lambda_2=0$，$\bm{\eta}_2 = \frac{1}{\sqrt{5}}(-1, 2)^T$.

        令$\bm{Q}_2 = (\bm{\eta}_1, \bm{\eta}_2) = \frac{1}{\sqrt{5}}\begin{bNiceMatrix} 2 & -1 \\ 1 & 2 \end{bNiceMatrix}$.

        由$\bm{Q}_1^T \bm{A} \bm{Q}_1 = \bm{\Lambda} = \bm{Q}_2^T \bm{B} \bm{Q}_2$，得$\bm{B} = \bm{Q}_2 \bm{Q}_1^T \bm{A} \bm{Q}_1 \bm{Q}_2^T = (\bm{Q}_1\bm{Q}_2^T)^T \bm{A} (\bm{Q}_1\bm{Q}_2^T)$.

        故取$\bm{Q} = \bm{Q}_1\bm{Q}_2^T = \dfrac{1}{5}\begin{bNiceMatrix} 1 & 2 \\ -2 & 1 \end{bNiceMatrix}\begin{bNiceMatrix} 2 & -1 \\ 1 & 2 \end{bNiceMatrix} = \begin{bNiceMatrix}[cell-space-top-limit=3pt] 0 & 1 \\ -1 & 1 \end{bNiceMatrix}$.
    \end{enumerate}
\end{solution}

\section{正定二次型与正定矩阵}

\begin{definition}[正定性]
    若对任意非零向量$x \neq 0$，都有$x^T \bm{A} x > 0$，则称二次型$f$为正定二次型，称对称矩阵$\bm{A}$为正定矩阵．
\end{definition}
\begin{definition}[顺序主子式]
    设 $\bm{A} = (a_{ij})_{n \times n}$，子式
    \[ D_k = \det
    \begin{bNiceMatrix}
        a_{11} & \cdots & a_{1k} \\
        \vdots & \ddots & \vdots \\
        a_{k1} & \cdots & a_{kk}
    \end{bNiceMatrix} \]
    称为矩阵 $\bm{A}$ 的 $k$ 阶顺序主子式 ($k = 1, 2, \cdots, n$)．
\end{definition}
\begin{theorem}[正定矩阵的判别]
    $n$阶实对称矩阵$\bm{A}$正定的充要条件是：
    \begin{enumerate}
        \item $\bm{A}$的所有特征值$\lambda_i > 0$．
        \item $\bm{A}$的正惯性指数$p = n$．
        \item $\bm{A}$合同于单位矩阵$\bm{E}$．
        \item $\bm{A}$的所有顺序主子式全大于0．
    \end{enumerate}
\end{theorem}


\begin{example}
    设实二次型$f(x_1, x_2, x_3) = (x_1-x_2+x_3)^2+(x_2+x_3)^2+(x_1+ax_3)^2$，其中$a$为实参数，求：
    \begin{tasks}[label=(\arabic*)](2)
        \task $f(x_1, x_2, x_3)=0$的解；
        \task $f(x_1, x_2, x_3)$的规范型；
        \task $f(x_1, x_2, x_3)$是否为正定二次型.
    \end{tasks}
\end{example}
\begin{solution}
    \begin{enumerate}[label = (\arabic*)]
        \item 由$f=0$且$f$为平方和形式，得
        \[ \begin{cases} x_1-x_2+x_3=0 \\ x_2+x_3=0 \\ x_1+ax_3=0 \end{cases} \qquad\qquad
        \begin{vNiceMatrix}[columns-width = 1.5em]
            1 & -1 & 1  \\
            0 & 1 & 1  \\
            1 & 0 & a 
            \end{vNiceMatrix} = a-2\]
        \begin{itemize}
            \item 当$a \neq 2$时，仅有零解$\bm{x}=\bm{0}$.
            \item 当$a = 2$时，有非零解$\bm{x} = k(-2, -1, 1)^T \; (k \neq 0)$.
        \end{itemize}
        \item 令$\bm{C} = \begin{bNiceMatrix}[columns-width = 1.2em] 1 & -1 & 1 \\ 0 & 1 & 1 \\ 1 & 0 & a \end{bNiceMatrix}$，$\det(\bm{C}) = a-2$.
        \begin{itemize}
            \item 当$a \neq 2$时，变换非退化，规范型为$y_1^2 + y_2^2 + y_3^2$.
            \item 当$a = 2$时，$\rank(\bm{C})=2$，配方法：
            \begin{align*}
                f &= (x_1 - x_2 + x_3)^2 + (x_2 + x_3)^2 + (x_1 + 2x_3)^2 \\
                &= 2(x_1 - \dfrac{1}{2}x_2 + \dfrac{3}{2}x_3)^2 + \dfrac{3}{2}(x_2 + \dfrac{2}{3}x_3)^2
            \end{align*}
            规范型为$y_1^2 + y_2^2$.
        \end{itemize}
        \item 正定充要条件为正惯性指数$p=3$.
        
        由(2)知，当$a \neq 2$时，$p=3$，是正定二次型；当$a=2$时，为半正定二次型.
    \end{enumerate}
\end{solution}
\section{习题}
\begin{exercise}
  已知$\bm{A}=\begin{bNiceMatrix}[columns-width = 1.5em]
    \lambda & 1 & 1 \\
    1 & \lambda & 1 \\
    1 & 1 & \lambda
    \end{bNiceMatrix}$，若$\lambda_1=3$是$\bm{A}$的一个单特征值，求\begin{enumerate}[label = (\arabic*)]
        \item 正交变换$\bm{x}=\bm{Q}\bm{y}$将二次型$f(x_1,x_2,x_3)=\bm{x}^T\bm{A}\bm{x}$化为标准形；
        \item $\bm{X}_{3 \times 2},$使得$\bm{X}^T(\bm{A}-3\bm{E})\bm{X}=\bm{0}$.
    \end{enumerate}
\end{exercise}
\begin{solution}

    (1) 令$|\bm{A}-k\bm{E}|=(\lambda - k + 2)(\lambda - k - 1)^2=0\implies$特征值为$\lambda+2, \lambda-1, \lambda-1$.

    已知$\lambda_1=3$是单特征值.

    若$\lambda-1=3$即$\lambda=4$，特征值为$6, 3, 3$，此时$3$是二重特征值，不合题意.

    若$\lambda+2=3$即$\lambda=1$，特征值为$3, 0, 0$，$3$是单特征值，符合题意.\\
    
    此时$\bm{A}=\begin{bNiceMatrix}[columns-width = 1.2em] 1 & 1 & 1 \\ 1 & 1 & 1 \\ 1 & 1 & 1 \end{bNiceMatrix}$.\\[0.5em]

    当特征值$\mu_1=3$时，解$(\bm{A}-3\bm{E})\bm{x}=\bm{0}$得基础解系$\bm{\xi}_1=(1,1,1)^T$，
    单位化得$\bm{e}_1=\frac{1}{\sqrt{3}}(1,1,1)^T$.

    当特征值$\mu_2=\mu_3=0$时，解$\bm{A}\bm{x}=\bm{0}$即$x_1+x_2+x_3=0$，
    可取基础解系$\bm{\xi}_2=(-1,1,0)^T,\bm{\xi}_3=(1,1,-2)^T$.
    正交化、单位化得$\bm{e}_2=\frac{1}{\sqrt{2}}(-1,1,0)^T, \bm{e}_3=\frac{1}{\sqrt{6}}(1,1,-2)^T$.

    令$\bm{Q}=(\bm{e}_1,\bm{e}_2,\bm{e}_3)$，
    则正交变换$\bm{x}=\bm{Q}\bm{y}$将二次型化为标准形$f=3y_1^2$.

    (2) 令$\bm{X}=(\bm{x}_1, \bm{x}_2),\,\bm{B}=\bm{A}-3\bm{E}$, 则
    \[ \bm{X}^T\bm{B}\bm{X} = \begin{bNiceMatrix} \bm{x}_1^T\bm{B}\bm{x}_1 & \bm{x}_1^T\bm{B}\bm{x}_2 \\ \bm{x}_2^T\bm{B}\bm{x}_1 & \bm{x}_2^T\bm{B}\bm{x}_2 \end{bNiceMatrix} = \bm{0} \]
    即需$\bm{x}_1, \bm{x}_2$满足$\bm{x}_i^T\bm{B}\bm{x}_j=0,\, i,j=1,2$.
    令$\bm{x} = \bm{Qy}$，则
    \[ \bm{x}^T\bm{B}\bm{x} = \bm{y}^T\bm{Q}^T\bm{B}\bm{Q}\bm{y} = \bm{y}^T\begin{bNiceMatrix} 0 & 0 & 0 \\ 0 & -3 & 0 \\ 0 & 0 & -3 \end{bNiceMatrix}\bm{y} =- 3y_2^2 - 3y_3^2 =0 \implies y_2 = y_3 = 0 \]
    即$\bm{x}$为$\bm{e}_1$的倍数，取$\bm{x}_1=k_1\sqrt{3}\bm{e}_1, \bm{x}_2=k_2\sqrt{3}\bm{e}_1, k_1,k_2 \in \mathbb{R}$.
    则
    \[\bm{X}=\begin{bNiceMatrix} k_1 & k_2 \\ k_1 & k_2 \\ k_1 & k_2 \end{bNiceMatrix} \]
\end{solution}

\begin{exercise}
    设$\bm{A}$为$3$阶实矩阵，$\bm{E}$为$3$阶单位矩阵，三维列向量$\bm{\alpha}=(1,2,3)^T$.\,现将$\bm{E}$的第一行与第二行互换得到矩阵$\bm{C}$，若$\bm{C}^T\bm{AC}=\bm{\alpha\alpha}^T+k\bm{E},\,$回答下列问题：
    \begin{enumerate}[label = (\arabic*)]
        \item 证明：$\bm{A}$为实对称矩阵；
        \item 求可逆矩阵$\bm{P}$，使得$\bm{P}^{-1}\bm{AP}$为对角矩阵；
        \item 若$f(x_1,x_2,x_3)=\bm{x}^T\bm{Ax}$为正定二次型，求$k$的取值范围.
    \end{enumerate}
\end{exercise}
\begin{solution}
    (1) 易知$\bm{C}$为正交矩阵且$\bm{C}=\bm{C}^T$. 令$\bm{B}=\bm{\alpha\alpha}^T+k\bm{E}$，显然$\bm{B}^T=\bm{B}$.
    由$\bm{C}^T\bm{AC}=\bm{B}$得$\bm{A}=\bm{C}\bm{B}\bm{C}^T$.
    $\bm{A}^T = (\bm{C}\bm{B}\bm{C}^T)^T = \bm{C}\bm{B}^T\bm{C}^T = \bm{C}\bm{B}\bm{C}^T = \bm{A}$，故$\bm{A}$为实对称矩阵.

    (2) 矩阵$\bm{A}$与$\bm{B}$相似，$\bm{B}$的特征值为$\bm{\alpha\alpha}^T$的特征值加$k$.
    $\bm{\alpha\alpha}^T$的特征值为$\text{tr}(\bm{\alpha\alpha}^T)=1^2+2^2+3^2=14$和$0,0$.
    故$\bm{A}$的特征值为$14+k, k, k$.
    
    易得$\bm{A}=\begin{bNiceMatrix}
        4 + k & 2 & 3 \\
        2 & 1 + k & 6 \\
        3 & 6 & 9 + k
    \end{bNiceMatrix}$

    当特征值$\lambda_1=14+k$时，解$(\bm{A}-(14+k)\bm{E})\bm{x}=\bm{0}$得基础解系$\bm{\xi}_1=(2,1,3)^T$，

    当特征值$\lambda_2=\lambda_3=k$时，解$(\bm{A}-k\bm{E})\bm{x}=\bm{0}$得基础解系$\bm{\xi}_2=(-1,2,0)^T, \bm{\xi}_3=(-3,0,2)^T$.

    取$\bm{P}=(\bm{\xi}_1, \bm{\xi}_2, \bm{\xi}_3)$，则$\bm{P}^{-1}\bm{AP} = \begin{bNiceMatrix}[columns-width = 1.5em]
        14+k & 0 & 0 \\
        0 & k & 0 \\
        0 & 0 & k
    \end{bNiceMatrix}$.


    (3) $f$为正定二次型 $\iff$ $\bm{A}$的特征值全大于0.
    即 $14+k>0$ 且 $k>0$，解得 $\boxed{k>0}$.
\end{solution}
\ifx\allfiles\undefined
\end{document}
\fi